\section*{Sets and Indices}

\[
T = \{1,\dots,n\}
\]
set of planning stages, indexed by \(t\).

For this instance,
\[
n=10.
\]

\section*{Parameters}

\[
r_t \qquad (t \in T)
\]
tool requirement in stage \(t\) (from \texttt{tool\_requirement}).

\[
c^{\text{new}}
\]
cost of buying one new tool (code data: \texttt{cost\_new\_tool}).

\[
c^{\text{fast}}
\]
cost of sending one tool to fast repair (code data: \texttt{cost\_fast\_repair}).

\[
c^{\text{slow}}
\]
cost of sending one tool to slow repair (code data: \texttt{cost\_slow\_repair}).

\[
\tau^{\text{fast}}
\]
fast repair duration in stages (code data: \texttt{fast\_repair\_duration}).

\[
\tau^{\text{slow}}
\]
slow repair duration in stages (code data: \texttt{slow\_repair\_duration}).

For the provided data,
\[
c^{\text{new}}=10,\quad
c^{\text{fast}}=3,\quad
c^{\text{slow}}=1,\quad
\tau^{\text{fast}}=1,\quad
\tau^{\text{slow}}=3,
\]
and
\[
(r_1,\dots,r_{10})=(3,5,2,4,6,5,4,3,2,1).
\]

\section*{Decision Variables}

For each \(t \in T\),

\[
x_t \in \mathbb{Z}_{\ge 0}
\qquad \text{(code: \texttt{x[t]})}
\]
number of new tools purchased in stage \(t\).

\[
y_t \in \mathbb{Z}_{\ge 0}
\qquad \text{(code: \texttt{y[t]})}
\]
number of tools sent to fast repair after stage \(t\).

\[
z_t \in \mathbb{Z}_{\ge 0}
\qquad \text{(code: \texttt{z[t]})}
\]
number of tools sent to slow repair after stage \(t\).

\[
c_t \in \mathbb{Z}_{\ge 0}
\qquad \text{(code: \texttt{c[t]})}
\]
number of clean/available tools carried at the end of stage \(t\), for \(t=0,\dots,n\).

\[
d_t \in \mathbb{Z}_{\ge 0}
\qquad \text{(code: \texttt{d[t]})}
\]
number of dirty tools waiting after stage \(t\), for \(t=0,\dots,n\).

\section*{Objective}

Minimize total purchasing and repair cost:
\[
\min \sum_{t=1}^n \left(c^{\text{new}} x_t + c^{\text{fast}} y_t + c^{\text{slow}} z_t\right).
\]

\section*{Constraints}

Initial conditions:
\[
c_0 = 0,
\]
\[
d_0 = 0.
\]

Stage-by-stage clean-tool balance, for all \(t \in T\):
\[
x_t + c_{t-1}
+ \mathbf{1}_{\{t-\tau^{\text{fast}}-1 \ge 1\}}\, y_{\,t-\tau^{\text{fast}}-1}
+ \mathbf{1}_{\{t-\tau^{\text{slow}}-1 \ge 1\}}\, z_{\,t-\tau^{\text{slow}}-1}
= r_t + c_t .
\]

This exactly matches the code logic that makes fast repairs sent in stage \(s\) available in stage \(s+\tau^{\text{fast}}+1\), and slow repairs sent in stage \(s\) available in stage \(s+\tau^{\text{slow}}+1\).

Dirty-tool balance, for all \(t \in T\):
\[
r_t + d_{t-1} = y_t + z_t + d_t .
\]

Integrality and nonnegativity:
\[
x_t,y_t,z_t \in \mathbb{Z}_{\ge 0} \qquad \forall t\in T,
\]
\[
c_t,d_t \in \mathbb{Z}_{\ge 0} \qquad \forall t=0,\dots,n.
\]

\section*{Notes on Implementation}

The implemented model is a mixed-integer linear program solved with Gurobi through PuLP.

The balance equations should be interpreted literally as written in the source code. In particular, a tool used in stage \(t\) enters the dirty stock after that stage, and if it is sent to repair in stage \(t\), it returns only in stage \(t+\tau+1\), not \(t+\tau\). Thus the code uses the terms
\[
y_{t-\tau^{\text{fast}}-1}, \qquad z_{t-\tau^{\text{slow}}-1},
\]
with zero contribution whenever the index is less than \(1\).

There is no explicit constraint forcing all dirty tools to be repaired by the end of the horizon; leftover dirty tools \(d_n\) are allowed. Likewise, end-of-horizon clean inventory \(c_n\) is allowed, although it is not rewarded. These modeling choices are part of the implemented formulation and are not tightened further here.
